
\subsubsection{Overview}

The proposed approach operates as a meta-reasoning layer that monitors proof search without modifying the base tactic suggestion system. The method extracts confidence trajectories from softmax distributions, computes stability metrics, and combines these with symbolic divergence indicators.

\subsubsection{Confidence Geometry Extraction}

For each proof step t, the LLM produces a softmax distribution p_t over candidate tactics. Shannon entropy is computed as:

H_t = -Σ_i p_{t,i} log p_{t,i}

The entropy trajectory {H_1, H_2, ..., H_T} over the first T steps is characterized by its variance:

σ²_H = (1/(T-1)) Σ_t (H_t - H̄)²

The design documents specify T=15 steps. The rationale is that variance captures trajectory instability: successful proofs should navigate smoothly through familiar state space (low variance), while divergent searches should show erratic confidence patterns (high variance).

The geometric interpretation is that entropy serves as an implicit distance metric from the training distribution manifold. On-manifold states yield stable entropy; off-manifold states yield erratic entropy.

\subsubsection{Symbolic Signal Extraction}

Two symbolic metrics are proposed:

\textbf{State Hash Collisions}: A hash table tracks encountered proof states. Collisions indicate revisited states, potentially signaling cyclic behavior.

\textbf{Proof State Growth}: The total size of active goals is tracked. Exponential growth flags potential syntactic explosion.

\subsubsection{Signal Combination}

The design documents describe multiple detector variants tested in ablation studies:

\begin{itemize}
\item Single signals: confidence-only, symbolic-only, search-only
\item Pairwise combinations: confidence+symbolic, confidence+search, symbolic+search
\item Hybrid: 3-signal voting with k=2-of-3 threshold
\end{itemize}

The pairwise confidence+symbolic detector uses OR logic:

flag = (σ²_H > θ_conf) OR (collisions > θ_sym)

Thresholds are set via median of the timeout distribution.

\subsubsection{Computational Overhead}

The design documents estimate approximately 15\% overhead based on complexity analysis. This has not been measured in actual deployment.
